\section{Limitations}
\label{sec:limitations}

\paragraph{Ground truth is accession-level and not fully independent.} Correctness
is judged against PDBe SIFTS and RCSB entity references. These are curated
independently of the software's alignment, but they are not independent of the
cross-references the accession check reads, so a result of zero false associations
at M4 shows consistency between two views of related curated data, not an
independent audit. A structure that is the right protein but catalogued under a
different accession is counted as a false association; the tiering of
Section~\ref{sec:validation-study} shows how much of the total this is, but the
headline rates should be read as accession-level.

\paragraph{The sample is not representative.} The \N{ds-n} proteins are a random
draw from reviewed UniProt entries, stratified by length and category, taken at one
time from one release. About \N{ds-no} of \N{ds-n} have no deposited structure,
which is typical of Swiss-Prot, and the rates are comparisons between policies on
this set, not estimates of prevalence in any user's workload.

\paragraph{Entries without a PDB-format file cannot be validated.} The RCSB
provider downloads PDB-format files, which do not exist for large or recent
entries. \N{gate-unmeasured} of the \N{gate-total} correct naive hits and
\N{n1-all-tier-unm} of the \N{n1-all-false} naive false associations could not be
validated for this reason (every case an HTTP 404). They cannot be accepted, which
costs sensitivity. Validating from the polymer-entity sequence when no PDB file
exists is a bounded fix we have not made.

\paragraph{M2 is not evidence of homology.} The related-record thresholds sit at the
alignment's noise floor: \N{nf-m2-pct}\% of unrelated pairs satisfy them
(Section~\ref{sec:threshold}). We did not retune them on the evaluation data.
Identity is computed over the aligned region only and is inflated for short
subject chains.

\paragraph{The strict policy has low sensitivity.} A correct structure was
returned for \N{s1-has-sens}\% of proteins that have one. Most of the loss is
legitimate (partial structures, no file to validate), but a user who needs
coverage must use the permissive operating point and accept its same-sequence and
near-identical false associations.

\paragraph{The adversarial benchmark is limited and partly unverified.} B5 has only
\N{adv-B5-n} scored cases. B1 depends on a Pfam grouping and a crude sequence-similarity
estimate used only for construction, and a pair that reaches M4 is most likely a
construction error that this benchmark cannot distinguish from a cross-reference
error. B3's label is an assumption, not ground truth, and is not reported as a
failure rate. \N{adv-unresolved} of \N{adv-total} built cases could not be scored,
so the scored cases are not a random subset. The frozen set was produced by a
script whose pair ordering was not reproducible, so regenerating it gives a different sample.

\paragraph{The 25-protein re-adjudication reproduces today's archive, not the
prototype's.} It cannot show what v1.0 retrieved in August 2026, only what a
sequence search returns now. It shows the prototype's discordance flags do not
demonstrate wrong structures; it does not show what the prototype did.

\paragraph{The live service benchmark is small and its M4 rate is a coverage figure.}
The 69-protein set was curated to include difficult cases and AlphaFold DB is
accession-keyed, so its M4 rate says nothing about validation. Service behaviour
varies with time; the dated runs retained in the repository document that variation.

\paragraph{Fault injection and usability.} The 80\% recovery figure is a property
of the injected scenarios, not an expected service-level rate. Interaction-step
counts are a proxy for effort; no human-subject usability study has been run, and
the protocol for one is included in the repository. The statement-coverage figure,
the standalone installer and the archival DOI are pending
(\todoinline{insert \texttt{pytest --cov} output, installer status and Zenodo DOI}).

\paragraph{Scope.} Structure retrieval covers a single query chain; multi-chain
assemblies are validated by their best-matching chain, and complexes in which the
queried protein is one of many chains are classified on that chain only.
